\documentclass[10pt, a4paper, landscape]{article}

\usepackage[T1]{fontenc}
\usepackage[utf8]{inputenc}
\usepackage{tgpagella}
\usepackage{tgcursor}
\usepackage{microtype}
\usepackage[spanish, es-noshorthands]{babel}
\usepackage[table, xcdraw]{xcolor}
\usepackage{booktabs}
\usepackage{tabularx}
\usepackage{array}
\usepackage{amsmath, amssymb}
\usepackage[most]{tcolorbox}
\usepackage[margin=1.4cm]{geometry}
\usepackage{fancyhdr}

\pagestyle{fancy}
\fancyhf{}
\setlength{\headheight}{16pt}
\lhead{\textbf{IMT-342}}
\chead{Bitácora Reproducible, Backlog W13 y Exit Ticket}
\rhead{Semana 12 - Sesión 2}
\lfoot{Compromisos Verificables hacia W13}
\rfoot{Página \thepage}
\renewcommand{\headrulewidth}{0.4pt}
\renewcommand{\footrulewidth}{0.4pt}

\definecolor{ucbblue}{RGB}{0, 51, 102}

\begin{document}

\begin{center}
    \Large\textbf{\textcolor{ucbblue}{Bitácora de Decisiones, Backlog hacia W13 (Máx. 3 Tareas) y Exit Ticket Individual}}[cite: 16]
\end{center}

\section*{1. Backlog de Pruebas Reproducibles hacia W13 (Máximo 3 Tareas por Equipo)[cite: 16]}
El producto del trabajo autónomo no es \textit{«investigar algo»}, sino producir una decisión reproducible respaldada por un archivo versionado, captura autorizada con origen, log fechado, tabla reproducible o enlace a commit concreto[cite: 16]. Si no existe acceso autorizado al ABB, toda prueba en hardware permanece \texttt{BLOCKED} y el equipo avanza en documentación comparativa, stack ROS local y ejecución de Gate 2 offline[cite: 16].

\begin{center}
\scriptsize
\setlength{\tabcolsep}{3.5pt}
\renewcommand{\arraystretch}{1.55}
% Suma de \hsize = 1.05 + 0.80 + 0.90 + 1.15 + 1.05 + 0.85 + 1.00 + 1.20 = 8.0 para las 8 columnas X (+ 1 col c)
\begin{tabularx}{\textwidth}{|c|>{\hsize=1.05\hsize\raggedright\arraybackslash}X|>{\hsize=0.80\hsize\raggedright\arraybackslash}X|>{\hsize=0.90\hsize\raggedright\arraybackslash}X|>{\hsize=1.15\hsize\raggedright\arraybackslash}X|>{\hsize=1.05\hsize\raggedright\arraybackslash}X|>{\hsize=0.85\hsize\centering\arraybackslash}X|>{\hsize=1.00\hsize\raggedright\arraybackslash}X|>{\hsize=1.20\hsize\raggedright\arraybackslash}X|}
\hline
\rowcolor{ucbblue!15}
\textbf{\#}[cite: 16] & \textbf{Hipótesis Técnica a Someter a Prueba}[cite: 16] & \textbf{Responsable y Suplente}[cite: 16] & \textbf{Dependencia / Permiso Requerido}[cite: 16] & \textbf{Procedimiento Permitido (Sin Mov. Físico)}[cite: 16] & \textbf{Evidencia Verificable (Archivo / Commit / Log)}[cite: 16] & \textbf{Criterio} (\texttt{PASS} / \texttt{FAIL} / \texttt{BLOCKED})[cite: 16] & \textbf{Incertidumbre Remanente}[cite: 16] & \textbf{Siguiente Decisión según Resultado}[cite: 16] \\ \hline
\rowcolor{gray!10}
\textbf{Ej.}[cite: 16] & Identificar el modelo exacto de controlador y asociarlo al manual de su versión[cite: 16]. & Coord. HW / Suplente SW[cite: 16] & Permiso de operador autorizado del lab.[cite: 16] & Inspección visual / foto de placa autorizada o solicitud formal si no hay acceso[cite: 16]. & Foto fechada + ficha actualizada + PDF manual (o cargo de solicitud)[cite: 16]. & \texttt{PASS} si hay placa y manual; \texttt{BLOCKED} sin acceso[cite: 16]. & Opciones de software internas aún no leídas[cite: 16]. & Habilitar cruce de versión con Ruta A/B en W13[cite: 16]. \\ \hline
\textbf{1}[cite: 16] & & & & & & & \\[0.6cm] \hline
\textbf{2}[cite: 16] & & & & & & & \\[0.6cm] \hline
\textbf{3}[cite: 16] & & & & & & & \\[0.6cm] \hline
\end{tabularx}
\end{center}

\vspace{0.2cm}
\section*{2. Exit Ticket Individual de Cierre (Segmento 6 — 10 min)[cite: 16]}
\begin{tcolorbox}[colback=white, colframe=ucbblue, title=\textbf{Verificación Individual de Salida — Estudiante:} \rule{7cm}{0.15mm} \textbf{Grupo:} \rule{2cm}{0.15mm} {[cite: 16]}, left=2mm, right=2mm, top=1.5mm, bottom=1.5mm]
    \scriptsize\raggedright
    \textbf{1. Indique un dato concreto del controlador ABB ya verificado con fuente, o el bloqueo exacto que impide conocerlo y quién debe autorizar el acceso[cite: 16]:}\\[0.35cm]
    \rule{\textwidth}{0.15mm}\\[0.25cm]
    \textbf{2. Describa una prueba de comunicación sin movimiento (Gate 1) y cuál es su observable exacto que distingue un estado real de una simulación[cite: 16]:}\\[0.35cm]
    \rule{\textwidth}{0.15mm}\\[0.25cm]
    \textbf{3. Enuncie una condición obligatoria de \texttt{NO-GO} vigente hoy antes de cualquier ensayo físico en el robot[cite: 16]:}\\[0.35cm]
    \rule{\textwidth}{0.15mm}\\[0.25cm]
    \textbf{4. Explique una razón geométrica por la cual dos modelos de Cinemática Directa (FK) correctos pueden diferir numéricamente debido a los marcos $\{B\}$ o $\{TCP\}$[cite: 16]:}\\[0.35cm]
    \rule{\textwidth}{0.15mm}
\end{tcolorbox}

\end{document}
